\chapter{Free energies and alloy thermodynamics}
\label{ch:freeenergy}

Three modules answer questions that a single energy cannot. Thermodynamic
integration produces an \emph{absolute} free energy from molecular dynamics;
the ECI fit turns a set of computed alloy configurations into a predictive
model; and the Cluster Variation Method turns that model into a configurational
entropy and, where the lattice allows it, an order--disorder temperature.

\section{Thermodynamic integration}
\label{sec:ti}

\menu{Simulation}\menusep\menu{Thermodynamic Integration\dots} computes $F$, and
$G = F + PV$, in absolute terms.

That is not something a simulation can measure. Energy, temperature, pressure
and stress are all mechanical averages over a trajectory; the free energy is the
logarithm of a partition function, and no average of any observable equals it.
Thermodynamic integration recovers it by building a \emph{reversible} path from
a reference system whose free energy is known in closed form:
\begin{equation}
  U(\lambda) = (1-\lambda)\,U_\mathrm{ref} + \lambda\,U_\mathrm{target},
  \qquad
  \frac{\partial F}{\partial\lambda}
    = \Bigl\langle \frac{\partial U}{\partial\lambda} \Bigr\rangle_\lambda
    = \bigl\langle U_\mathrm{target} - U_\mathrm{ref} \bigr\rangle_\lambda ,
  \label{eq:ti-coupling}
\end{equation}
\begin{equation}
  \Delta F = \int_0^1
    \bigl\langle U_\mathrm{target} - U_\mathrm{ref} \bigr\rangle_\lambda
    \,\mathrm{d}\lambda ,
  \qquad
  F = F_\mathrm{ref} + \Delta F ,
  \qquad
  G = F + PV .
  \label{eq:ti-integral}
\end{equation}
Each $\lambda$ is one MD run --- a \emph{window} --- reporting one number. Every
energy is per simulation cell, matching the phonon thermodynamics of
Chapter~\ref{ch:elph}'s neighbour module; per-atom values are reported beside
them with the convention spelled out in the field names.

\subsection{The reference systems}

The method rests entirely on $F_\mathrm{ref}$ being \emph{exact}. An approximate
reference does not add noise; it adds a bias no amount of MD removes.

\begin{description}
  \item[\ui{Ideal gas (Sackur-Tetrode)}] $F_\mathrm{id} = k_\mathrm{B}T
    \sum_s N_s [\ln(\rho_s\Lambda_s^3) - 1]$, summed over \emph{species}: atoms
    are grouped by mass, and each group carries its own thermal wavelength and
    its own $N!$. The factorial is the point --- dropping it gives a number that
    scales almost right and is wrong by $Nk_\mathrm{B}T(\ln N - 1)$. Exact, with
    no caveat, and the natural reference for a liquid. It is also the reference
    with the worst endpoint behaviour (\S\ref{sec:ti-endpoint}).
  \item[\ui{Einstein crystal (Frenkel-Ladd)}]
    $F_\mathrm{Ein} = 3Nk_\mathrm{B}T\ln(\beta\hbar\omega)$ with
    $\omega = \sqrt{\alpha/m}$. Exact, and free of the endpoint singularity
    because tethered atoms never overlap. There is deliberately \emph{no}
    $1/N!$ here --- the opposite of the ideal-gas case, and exactly where a
    copy-paste between the two goes wrong: Einstein atoms are tethered to
    distinguishable sites, cannot exchange, and a crystal samples one
    permutation. The \ui{Hold the centre of mass fixed} box drives \emph{both}
    the \texttt{FixCom} constraint in the run and the matching $O(\ln N)$
    finite-size correction in the closed form, so the two cannot disagree.
  \item[\ui{Lennard-Jones fluid}] See the warning below.
\end{description}

\begin{calwarn}
\textbf{The Lennard-Jones reference cannot carry a liquid.} There is no exact
closed form for the free energy of an LJ fluid. Calango evaluates the exact
Hirschfelder--Curtiss--Bird second-virial series, $\beta A_\mathrm{ex}/N =
B_2\rho + O(\rho^2)$ --- an exact expansion, not a fit --- which is quantitative
only for a reduced density $\rho^* = (N/V)\sigma^3 \le 0.05$. Above that
the reference is marked invalid and no free energy is reported at all.

A liquid is an order of magnitude denser than that, so \textbf{for an actual
liquid use the ideal gas or the Einstein crystal}. The LJ entry exists to put an
LJ fluid \emph{on} a path, as the target of its own ideal-gas integration whose
excess free energy you then supply back --- not to start one at liquid density.
The alternative would have been to transcribe a 33-parameter equation-of-state
fit; an unverifiable table of coefficients that silently produces plausible
numbers is worse than a function that refuses.
\end{calwarn}

\subsection{The endpoint singularity}
\label{sec:ti-endpoint}

This is the failure mode of the method. With linear coupling
\eqref{eq:ti-coupling} and an ideal-gas reference, atoms at small $\lambda$ are
almost free and may sit on top of one another; a target with a hard repulsive
core then charges an unbounded energy for those configurations, and
$\langle\partial U/\partial\lambda\rangle$ diverges as $\lambda \to 0$. A
uniform grid integrates that divergence without noticing: every window returns a
finite number, the quadrature returns a finite number, and the answer is wrong.

Three layers stand against it.

\begin{enumerate}
  \item \textbf{The default schedule never samples $\lambda = 0$.}
    Gauss--Legendre nodes are strictly interior to $[0,1]$ and already cluster
    towards both ends, which is where the integrand bends.
  \item \textbf{A runtime detector.} The variance and the magnitude of the two
    end windows are compared against the \emph{median over the path} --- a
    ratio, not a fixed energy, because the integrand's scale is whatever the
    cohesive energy happens to be and only its \emph{shape} says whether a
    divergence is being integrated. Above $8\times$ either way, the result
    carries the diagnosis and both ratios.
  \item \textbf{A warning in the wizard.} Selecting the ideal gas together with
    a schedule whose first node is $\lambda = 0$ (Uniform, or Power law) turns a
    red note on the settings page, before anything has been paid for.
\end{enumerate}

\begin{calnote}
\textbf{There is no soft-core coupling, and there cannot be one here.} The
textbook fix replaces the linear mixing with a coupling that regularises the
\emph{pair form} of the potential. Calango's target is an arbitrary ASE
calculator --- MACE, GPAW, xTB, a LAMMPS force field --- and none of them
exposes a pair form to modify. There is no general way to soften a
machine-learning potential or a self-consistent DFT energy. So the module avoids
the endpoint, detects it when it bites anyway, and says so; it does not pretend
to have solved it.
\end{calnote}

\subsection{Settings}

\begin{description}
  \item[\ui{$\lambda$ schedule}] Gauss--Legendre (interior nodes) by default;
    also Uniform, Power law (clustered at $\lambda = 0$) and Clustered at both
    ends.
  \item[\ui{Quadrature}] Gauss--Legendre, Simpson or trapezoid. Gauss weights on
    non-Gauss nodes are not an approximation of anything, so that combination is
    \emph{refused} rather than silently reweighted; Simpson needs a uniform grid
    and falls back to the trapezoid, saying so. The two combos are interlocked
    in the wizard for the same reason.
  \item[\ui{$\lambda$ windows}] Default 12. Each is an independent MD run, and
    this is also the knob that fixes the quadrature error.
  \item[\ui{Dispatch as separate jobs}] Default 1. See the warning below.
  \item[\ui{Also sweep $\lambda$ backwards}] The hysteresis check. A reversible
    path must give the same $\Delta F$ either way; it will not if the windows
    were under-equilibrated, if the system changed phase along the path, or if
    the coupling drove it through a barrier --- and none of those looks like
    noise. Doubles the cost and needs one sequential chain, so it is withdrawn
    when the run is split.
  \item[\ui{Ensemble}] NVT--Langevin by default. \textbf{NVE is deliberately
    absent}: $\langle\partial U/\partial\lambda\rangle$ must be a canonical
    average at \emph{one} temperature, and a microcanonical window sits at a
    different temperature for every $\lambda$. NPT is offered but warned about
    --- the reference free energy is a function of $V$, and under a barostat
    there is no single volume to evaluate it at.
  \item[\ui{Equilibration steps / window}] Default 2000, and not optional in
    practice. Averaging over the transient biases every window in the
    \emph{same} direction, so the bias survives the $\lambda$ integral instead
    of cancelling --- and it looks nothing like noise on a plot.
  \item[\ui{Production steps / window}, \ui{Sampling interval}] 10\,000 and 10.
\end{description}

\begin{calwarn}
\textbf{$\lambda$ windows run sequentially.} Splitting the path into jobs
produces contiguous slices that go through the ordinary job queue, which has a
single runner, so they execute one after another exactly as the unsplit path
would. What splitting buys is \emph{resumability} --- a dead slice costs its own
windows, not the path --- and \emph{cluster dispatch}. It is not a wall-time
speed-up, and should not be budgeted as one.
\end{calwarn}

\subsection{Two error bars, and which one to spend on}

The generated script writes the \textbf{raw} $\partial U/\partial\lambda$
series, not merely its mean and variance: the autocorrelation time cannot be
recovered from moments, and MD samples are not independent. Calango computes the
integrated autocorrelation time by \textbf{Sokal's automatic windowing} ---
truncated, because the tail of the correlogram is pure noise whose variance
grows with the number of terms added, so the naive full sum does not merely lose
precision, it does not converge --- and reports
\begin{equation}
  \sigma_\mathrm{mean}
    = \sqrt{\frac{\mathrm{var}\cdot\tau_\mathrm{int}}{N}} .
\end{equation}
A naive $\sigma/\sqrt{N}$ over correlated samples under-reports by
$\sqrt{2\tau}$, routinely a factor of three: exactly the size of the
discrepancies people then explain away. A block-averaging estimate is computed
independently as a cross-check, and the two disagreeing by more than $2.5\times$
is reported as a production run too short for either.

Because every rule here is linear, the quadrature is expressed as
\emph{weights} rather than as a number, which lets the statistical error
propagate exactly:
\begin{equation}
  \sigma_I = \sqrt{\textstyle\sum_i w_i^2 \sigma_i^2} .
\end{equation}
Separately, a \textbf{$\lambda$-grid error} is estimated by comparing the chosen
rule against a closed trapezoid on the same nodes and on every second node. When
it exceeds $3\times$ the statistical error the report says so in as many words:
\emph{the lambda grid, not the sampling, dominates the error --- add windows
rather than MD steps.}

\subsection{Partial failure}

The number of windows the run was \emph{supposed} to have is passed to the
integrator separately from the list that came back; passing the length of the
list defeats the check, which is the whole reason the parameter exists.

A path with a dead window is not a noisier path, it is a \emph{different
integral}: quadrature weights are a property of the node set, so dropping a node
and reweighting the survivors integrates a curve that was never sampled. A run
missing any window therefore reports
\begin{center}
\texttt{*** INCOMPLETE --- NO FREE ENERGY IS REPORTED ***}
\end{center}
and leaves $\Delta F$ at zero. The summary distinguishes \emph{missing} (no file
--- the window never ran) from \emph{failed} (a file with a failure status ---
it ran and died), which is what tells you whether to re-run a slice or to change
the physics.

\subsection{The integrand viewer}

The free energy is a \emph{number}; what goes wrong is a \emph{shape}. Every
characteristic TI failure is visible in the integrand and invisible in
$\Delta F$, so the plot is not decoration. It shows
$\langle\partial U/\partial\lambda\rangle$ against $\lambda$ with

\begin{itemize}[leftmargin=1.4em]
  \item $1\sigma$ error bars from the \emph{autocorrelation-corrected} estimate,
    not the raw variance --- which understates a correlated series by exactly
    the factor this module exists to compute;
  \item the shaded integral, because $\Delta F$ \emph{is} that area;
  \item failed windows as vertical rules, so a gap is visible rather than
    inferred from a warning list;
  \item the endpoint band, when the detector fired;
  \item forward and backward sweeps overlaid when the hysteresis check ran.
\end{itemize}

It opens \emph{by itself} when a run finishes. There is no menu entry for it:
an ``open the results'' action whose only job is to ask which file you meant is
a question the application can already answer. The other three ways in all skip
that question --- the \ui{Load Results\dots} button on the wizard's
\emph{Integration Path \& Sampling} stage, the same button on this window
(which walks from one run to the next), and the Processes panel, which offers
\ui{Thermodynamic Integration Viewer} on any job directory holding a
\texttt{ti.json}. \ui{Export Results\dots} writes the per-window table as
CSV --- with the run's conditions and the assembled $F$, $PV$ and $G$ as
commented header lines --- or the report text as shown. The quadrature, the statistics,
the closed-form references and the endpoint diagnostics all run \emph{at read
time} in C++, so a run downloaded from a cluster is analysed by the same code as
a local one.

\subsection{A worked example}

Liquid copper, 32 atoms at 1600\,K, MACE as the target, ideal-gas reference,
five Gauss--Legendre windows. This was a development run rather than a shipped
test, so the numbers are reported and not asserted.

\begin{center}
\small
\begin{tabular}{@{}lp{7.6cm}@{}}
\toprule
Quantity & Value \\
\midrule
$\langle\partial U/\partial\lambda\rangle$ across the path & $-53.6 \to -120.9$ eV \\
Per-sample variance across the path & $301 \to 1.17$ --- a factor of \textbf{257} \\
Endpoint detector & \textbf{fired}, at $\times 107$ against the path median \\
$\Delta F$ & $-106.35$ eV \\
Statistical error & 0.46 eV \\
$\lambda$-grid error & \textbf{6.57 eV} \\
\bottomrule
\end{tabular}
\end{center}

The conclusion is the opposite of the instinct: \textbf{more MD time would buy
nothing.} The sampling error is already 0.46\,eV and the grid error is
6.57\,eV --- $14\times$ larger. Five nodes cannot resolve an integrand whose
variance changes by a factor of 257 along the path. The run needs more windows,
clustered harder towards $\lambda = 0$; longer production runs would shrink the
smaller of the two errors and leave the answer exactly as wrong. This is
precisely why the two errors are reported separately rather than added.

\begin{calwarn}
\textbf{Only EMT and MACE have been exercised end to end.} GPAW and xTB targets
are supported by construction --- the target is an ordinary ASE calculator and
the coupling is engine-agnostic --- but they are untested, and a DFT target pays
a full SCF per MD step, per window. The C++ half is covered by the
\texttt{thermodynamic\_integration} test (161 assertions): the closed-form
references against their analytic values, the Gauss--Legendre rule against the
polynomial exactness that defines it, the autocorrelation estimator against
block averaging, and the refusal paths.
\end{calwarn}

\section{Effective cluster interactions}
\label{sec:eci}

\menu{Modules}\menusep\menu{Alloys}\menusep\menu{Effective Cluster Interactions
(ECI Fit)\dots} turns a completed Cluster Expansion run into a predictive model,
\begin{equation}
  E(\sigma)/\text{atom}
    = J_0 + \sum_\alpha m_\alpha J_\alpha \Phi_\alpha(\sigma) ,
\end{equation}
by regressing the computed energies on the cluster correlations
$\Phi_\alpha$ that the builder emitted alongside them.

\subsection{Why a plain least-squares solve is a trap}

A cluster-expansion design matrix is ill-conditioned \emph{by construction}:
orbits at nearly the same radius have nearly the same correlation on every
configuration one can afford to compute, so their columns are near-parallel; the
empty cluster is exactly parallel to the intercept; and the number of candidate
orbits grows faster with the cutoff than the number of configurations anyone is
willing to run, so the system is frequently underdetermined outright. Forming
$X^{\mathsf T}X$ squares the condition number of an already borderline matrix
and then inverts it. The failure is not a crash --- it is a fit with a perfect
training residual and ECIs of alternating sign and absurd magnitude that
predicts nothing. Every solver here goes through an SVD or through coordinate
descent, never through an explicit inverse.

\subsection{Methods}

\begin{description}
  \item[\ui{LASSO (selects clusters)}] $L_1$. \emph{Selects} orbits, driving
    them to exactly zero --- bit for bit, not merely small. The default, because
    deciding which clusters carry the physics is the central difficulty of a
    cluster expansion and this is the method that decides.
  \item[\ui{Ridge (keeps all clusters)}] $L_2$. Keeps every orbit and only
    shrinks it; never sparse. It cannot tell you that a cluster does not matter,
    only that its coefficient is small, which is a different statement.
  \item[\ui{ARD / Bayesian}] Learns one precision per orbit by evidence
    maximisation and prunes outright, with no regularisation parameter to
    choose. Sharper support recovery than lasso when it works; slower, and it
    can prune too hard on very small data sets.
\end{description}

Columns are centred and scaled before fitting and unscaled afterwards; without
that, neither penalty being scale-invariant, the \emph{units} of an orbit would
decide how hard it is penalised.

\subsection{The CV score is the diagnostic}

\begin{calwarn}
A cluster expansion is \textbf{not} judged by its training RMSE but by its
\textbf{cross-validation score}, the RMS error on configurations the fit never
saw. The two numbers moving in opposite directions as the regularisation is
relaxed \emph{is} the definition of overfitting, so both are shown side by side.
A fit whose CV score exceeds $3\times$ the training error is flagged in red:
\emph{this expansion fits its own data and does not predict.}
\end{calwarn}

\ui{CV folds} defaults to leave-one-out, which is the cluster-expansion
convention --- ``the CV score'' in the CE literature \emph{means} the
leave-one-out score. The fitted orbits are ranked by $m\cdot J$ rather than by
$J$: a large ECI on a multiplicity-1 orbit and a small one on a
multiplicity-24 orbit can be the same physics.

A run made \emph{before} correlations were emitted has energies and no design
matrix, and is \textbf{refused with instructions} rather than fitted against
whatever else is present --- the alternative being an ECI file of zeros that the
CVM downstream would turn into a plausible and entirely wrong entropy curve.

The builder now defaults the \ui{Triplet cutoff} to 3.0\,\AA{} (on), because a
pair-only basis is symmetric under A$\leftrightarrow$B at complementary
compositions and cannot distinguish A$_3$B from AB$_3$ --- an ensemble built
without triplets silently discards the term that chooses the ordered structure.
\ui{Quadruplet cutoff} stays off: each orbit adds $K^4$ histogram columns, and
with the usual handful of configurations the fit runs out of samples before it
runs out of clusters.

\subsection{Send to CVM}

\ui{Send to CVM\dots} opens the CVM module with the fitted nearest-neighbour
pair ECI already converted and applied.

\begin{calwarn}
In the $\pm 1$ correlation basis the pair energy is $J s_i s_j$ with $s = +1$
for A and $-1$ for B, so
\begin{equation}
  e_{AB} = -J_2 , \qquad e_{AA} = e_{BB} = +J_2
  \qquad\Longrightarrow\qquad J_2 > 0 \ \text{ORDERS} .
  \label{eq:eci-sign}
\end{equation}
Half the literature writes the opposite convention. Getting it backwards turns
an ordering alloy into a clustering one and \emph{nothing fails} --- you get a
smooth, plausible entropy curve for the wrong material. That is why the
conversion is done by shared code and why the receiving window shows the
provenance of the number it was given.
\end{calwarn}

\section{The Cluster Variation Method}
\label{sec:cvm}

\menu{Modules}\menusep\menu{Alloys}\menusep\menu{CVM / Alloy
Thermodynamics\dots} computes the configurational entropy of a substitutional
alloy, and --- with the long-range-order option --- its order--disorder
temperature.

The question is how much of an alloy's entropy is really
$k_\mathrm{B}\ln(\text{arrangements})$. The textbook
$S_\mathrm{ideal} = -k_\mathrm{B}\sum_i x_i\ln x_i$ assumes the species are
placed \emph{independently} on every site. They are not: if unlike neighbours are
favoured the alloy orders, if like neighbours are favoured it clusters, and
either way there are fewer distinguishable arrangements than the ideal count.
$S_\mathrm{ideal}$ is an \emph{upper bound} --- and for a high-entropy alloy
advertised as stabilised by configurational entropy, the gap between the bound
($\ln 5 = 1.609\,k_\mathrm{B}$ for an equiatomic quinary) and the truth is the
entire question.

Kikuchi writes the entropy as a sum over clusters with alternating-sign
corrections, so that correlations already accounted for by a cluster are not
counted again by its subclusters:
\begin{equation}
  S/k_\mathrm{B} = -\sum_\alpha a_\alpha
    \sum_{\text{configs}} \rho_\alpha \ln \rho_\alpha ,
\end{equation}
with $a_\alpha$ the Kikuchi--Barker coefficients, fixed by requiring every
subcluster's contribution to cancel down to exactly one net counting.

\subsection{A hierarchy, not a menu}

\begin{center}
\small
\begin{tabular}{@{}llp{6.4cm}@{}}
\toprule
Approximation & Also known as & What it sees \\
\midrule
Point       & Bragg--Williams          & Nothing. Sites independent; this
                                         \emph{is} $S_\mathrm{ideal}$, and no
                                         interaction enters. \\
Pair        & Bethe--Peierls--Guggenheim & Nearest-neighbour pair correlations.
                                         \textbf{Exact on a 1D chain.} \\
Tetrahedron & Kikuchi                  & Points, pairs, triangles and
                                         tetrahedra --- the smallest cluster
                                         containing the frustrated triangles FCC
                                         is built from. \\
\bottomrule
\end{tabular}
\end{center}

Each is the previous one plus more correlation, so \textbf{computing all three
is the deliverable}: the spread between them \emph{is} the size of the
correlation correction, and a single number from any one of them does not say
whether it mattered. For an FCC binary at $x = 0.5$ with
$e_{AB} = -0.02$\,eV, at 500\,K:
\begin{equation}
  S_\mathrm{pair} = 0.5358 \;<\;
  S_\mathrm{tet}  = 0.6220 \;<\;
  S_\mathrm{ideal}= 0.6931 \quad k_\mathrm{B} .
\end{equation}
The ordering is the physically meaningful part. The pair approximation cannot
see the frustration of the close-packed tetrahedron --- on FCC one cannot make
every nearest-neighbour triangle unlike-decorated --- so it believes the alloy
orders more than it does, and therefore \emph{over}estimates ordering.

\begin{caltip}
\textbf{Warren--Cowley $\alpha$} is reported at every temperature, and it is
what makes the temperature dependence legible: $S$ alone does not say
\emph{which way} the alloy departs from random. $\alpha = 1 - P(j|i)/x_j$ is
zero for a random alloy, negative for ordering and positive for clustering. It
is also the more sensitive probe --- the residual short-range order is
\emph{first} order in $\beta J$ while the entropy deviation is \emph{second}
order --- and in the disordered phase it is directly comparable against a
diffuse-scattering measurement, which the entropy is not.
\end{caltip}

\subsection{Triplets}

The tetrahedron approximation consumes a nearest-neighbour \emph{triplet}
interaction as well as a pair one, and this changes what the model can express
at all. A binary with pair interactions only is exactly symmetric under
A$\leftrightarrow$B at complementary compositions, so it assigns $x = 0.25$ and
$x = 0.75$ the same energy and cannot distinguish A$_3$B from AB$_3$. On an FCC
tetrahedron at 600\,K, $E(x{=}0.25) - E(x{=}0.75)$ measures
$9.4\times10^{-17}$\,eV pair-only --- machine zero, an identity rather than a
small residual --- and $2.41\times10^{-2}$\,eV with a triplet ECI.

A triangle is not a subcluster of a pair, so a pair model has nowhere to put
one. Triplets supplied with the Pair approximation are \emph{dropped}, with the
warning that the numbers below describe a pair-only model and must not be read
as the cluster expansion that was fitted. The window itself exposes a single
$e_{AB}$; triplets reach the solver through the Orchestration canvas, whose CVM
node supplies them only when the approximation is the tetrahedron.

\subsection{Long-range order and $T_\mathrm{c}$}

\begin{calwarn}
\textbf{The homogeneous solver cannot produce an order--disorder transition at
all.} One sublattice means every site is statistically identical and there is no
symmetry left to break. Run Cu$_3$Au through it and you get a perfectly smooth
$S(T)$ and $\alpha(T)$ straight through 663\,K --- not a bug, but the
approximation stating its own limits. For a transition, tick \ui{Allow
long-range order (4 sublattices)}.
\end{calwarn}

The FCC nearest-neighbour tetrahedron has exactly one vertex on each of the four
simple-cubic sublattices FCC decomposes into, so the tetrahedron distribution
can be indexed by \emph{position} = sublattice and \emph{value} = species ---
and, unlike the homogeneous case, it is not symmetric under permuting the
positions. That asymmetry \emph{is} the long-range order. Only the overall
composition is constrained; the per-sublattice compositions are free, which is
what lets L1$_2$ form at $x = 1/4$ with three A-rich sublattices and one B-rich
one.

Both branches are solved at every temperature and $T_\mathrm{c}$ is located by
\textbf{which has the lower free energy} --- by the free energies crossing, not
by an order parameter decaying to zero. The FCC L1$_2$ transition is
\textbf{first order} and $\eta$ \emph{jumps}; a search for $\eta \to 0$ would
either find nothing or find the spinodal, which is above $T_\mathrm{c}$ and is
not the transition.

Ticking the box also snaps the composition to the structure's stoichiometry
(L1$_2 \leftrightarrow x = 1/4$, L1$_0 \leftrightarrow x = 1/2$) --- without
which asking for L1$_2$ at $x = 0.5$ correctly reports ``no transition'' and
reads as a broken feature --- and raises the minimum temperature, for the reason
in \S\ref{sec:cvm-limits}.

\paragraph{Cu$_3$Au.} The tetrahedron CVM gives $k_\mathrm{B}T_\mathrm{c} =
1.9245\,V_2$ for L1$_2$ at $x = 1/4$, against $3.2806\,V_2$ from
Bragg--Williams on the same system; $V_2 = 29.69$\,meV reproduces the measured
663\,K. The mean-field comparison is a \emph{provable} check rather than a
tolerance: for L1$_0$ at $x = 1/2$ the Bragg--Williams transition has the closed
form $k_\mathrm{B}T_\mathrm{c} = 4V_2$ exactly --- derived by linearising the
mean-field equations, and equivalently by Fourier-transforming the FCC
nearest-neighbour interaction, whose minimum sits at the X point with
$J(X) = -4V_2$ --- which the solver returns as $3.99972$, and one-sidedly: the
mean-field L1$_0$ transition is \emph{second} order, so both the order parameter
and the free-energy difference vanish continuously at $T_\mathrm{c}$ and the
finite thresholds the bisection needs bite a hair below the true temperature.
CVM must come out \emph{below} mean field because it includes correlations mean
field ignores, and FCC is frustrated enough that it is not merely a little
below.

\begin{calnote}
The bracket \textbf{MC $<$ CVM $<$ mean field} is the robust statement, and both
its ends are provable. The particular literature values often quoted for the
middle --- about 1.89 for the tetrahedron CVM, about 1.74 for Monte Carlo, and
the 25--30\,meV window for the Cu--Au nearest-neighbour $V_2$ --- are recalled
from standard literature and are \emph{not} traced to a cited source anywhere in
this codebase. Calango measures 1.893 for L1$_0$ at $x = 1/2$ and 1.9245 for
L1$_2$ at $x = 1/4$; treat the agreement as encouraging rather than as
validation against a reference you can look up from here.
\end{calnote}

\subsection{What this is not}
\label{sec:cvm-limits}

\begin{calwarn}
\textbf{This is not a phase diagram.} Neither solver performs a common-tangent
construction, and both work at \emph{fixed composition}, comparing homogeneous
phases. There are therefore no two-phase fields, no solvus lines, and no
composition range over which an ordered phase is stable.
\end{calwarn}

\begin{itemize}[leftmargin=1.4em]
  \item \textbf{Nearest-neighbour interactions only.} No second-neighbour terms,
    which on FCC are what select between competing ordered structures in several
    systems. The four-sublattice solver additionally takes pair interactions
    only.
  \item \textbf{BCC borrows FCC's tetrahedron counting and is not
    quantitative.} The Kikuchi--Barker coefficients depend on how many pairs,
    triangles and tetrahedra share a site, which is a property of the lattice
    and nothing else; BCC needs the \emph{irregular} tetrahedron built from
    first and second neighbours. Choosing BCC does not fail loudly --- it
    produces a smooth, plausible, wrong entropy --- so the run says so in its
    warnings.
  \item \textbf{Configurational entropy only}; no vibrational, electronic or
    magnetic contribution. Long-range order is FCC-only (L1$_2$, L1$_0$).
\end{itemize}

Two pathologies of the truncated Kikuchi expansion on frustrated FCC are
reported rather than hidden, and neither is a bug in the iteration.

\begin{enumerate}
  \item \textbf{Negative entropy on the disordered branch.} Well below the
    transition the \emph{disordered} branch returns $S < 0$ --- at $x = 1/2$ it
    tends to $-0.929\,k_\mathrm{B}$. The Kikuchi expansion is not a probability
    and nothing forces its truncation to stay positive on a lattice this
    frustrated. It never affects the stable branch or $T_\mathrm{c}$, which is
    what you read; it does mean the disordered free energy is meaningless deep
    inside the ordered field.
  \item \textbf{A domain-mixed fixed point.} At $x = 1/2$, below roughly
    $0.35\,T_\mathrm{c}$, the L1$_0$ start converges onto an equal mixture of
    the two symmetry-related L1$_0$ domains: exactly the ground-state energy,
    exactly $S = \ln 2$, and \emph{zero} long-range order, because the point
    marginals of the two domains average away while their pair correlations do
    not. It is a fixed point only because zeros in the pair marginals are
    absorbing under a multiplicative update, and a real crystal would pay a
    domain wall. Solutions with vanishing order parameter are therefore never
    accepted as ordered, and hitting one is reported --- which is why an
    L1$_0$ transition must be bracketed from inside the ordered field, and why
    ticking the checkbox raises the minimum temperature. L1$_2$ at $x = 1/4$
    does not suffer from this at any temperature.
\end{enumerate}

\subsection{Validation}

The tests check identities, not previous outputs. The pair approximation is
\emph{exact} on a 1D chain, and the 1D Ising chain has a closed-form
transfer-matrix solution: Calango matches it to $10^{-9}\,k_\mathrm{B}$ in
entropy and $10^{-12}$\,eV in energy, measuring residuals of order $10^{-16}$.
Both signs of $J$ are run, because a sign error in the energy convention passes
one and fails the other.

The zero-interaction limit pins the FCC Kikuchi--Barker coefficients
$(a_4, a_3, a_2, a_1) = (2, 0, -6, 5)$ --- the triangles genuinely vanish ---
together with the composition multiplier, by requiring $S$ to be \emph{exactly}
$S_\mathrm{ideal}$ when nothing interacts. It is checked at $x = 0.3$, not
$x = 0.5$, and measures $1.7\times10^{-14}\,k_\mathrm{B}$. The composition is
the point: an earlier derivation dropped the composition multiplier, giving
$w \sim (\prod x)^{7/8}$, which at $x = 0.5$ is completely invisible --- every
tuple picks up the same factor and normalisation hides it --- and wrong
everywhere else. The four-sublattice version of the same identity runs at
$x = 0.25$ and measures $2.2\times10^{-16}\,k_\mathrm{B}$.
